\section{ResearchPipeline: Dependency-Ordered Workers, Version-Bump Retries, and LaTeX Generation}
\label{sec:research}

\subsection{Overview}

The ResearchPipeline is UMAF's most architecturally sophisticated pipeline, implementing a 4-node LangGraph workflow \cite{d3mas,paperorchestrator} (\texttt{head $\rightarrow$ workers $\rightarrow$ reviewer $\rightarrow$ writer}) that transforms unstructured research topics into scored, ranked, LaTeX-formatted proposals. Three interlocking resilience mechanisms address the fundamental reliability problem in autonomous LLM research agents: dependency-ordered parallel execution with stop-on-failure, version-bump retry with checkpoint-based context reuse, and honest success reporting via \texttt{os.path.isfile()}.

\subsection{Dynamic Decomposition with Complexity-Scaled Sizing}

The head agent uses a three-tier fallback chain:

\textbf{Tier 1 --- LLM Decomposition}: The decomposer inherits from \texttt{BaseDecomposerRole} and overrides five template methods. The \texttt{\_sizing\_guide()} provides complexity heuristics: narrow topics $\rightarrow$ 3--5 sub-topics, moderate $\rightarrow$ 5--8, broad $\rightarrow$ 8--12.

\textbf{Tier 2 --- Timeout Protection}: The head node wraps the decomposition in a \texttt{ThreadPoolExecutor} with \texttt{HEAD\_TIMEOUT=300}s.

\textbf{Tier 3 --- Programmatic Fallback}: Splits on commas, ``and'', ``vs'', semicolons; creates a deep-dive subtopic per keyword; appends synthesis sub-topics; guarantees $\ge$2 sub-topics.

\subsection{Dependency-Ordered Execution via Topological Levels}

\texttt{BasePipeline.\_topological\_levels()} transforms a flat task list with dependencies into execution levels using a modified Kahn's algorithm---a pattern inspired by the MapReduce execution model \cite{mapreduce} and extended by agentic MapReduce approaches \cite{amapreduce,llmmapreduce}: key resolution (dependencies can be int IDs or string names), level construction (all independent tasks grouped), cycle detection and breaking (iterative edge removal prioritizing last dependencies), and fallback to parallel execution for unbreakable cycles.

\texttt{\_run\_workers\_with\_deps()} executes levels sequentially, injects \texttt{\_dependency\_outputs} into dependent tasks, and implements stop-on-failure: if any task fails and downstream levels remain, execution breaks immediately, deferring downstream dependents for retry.

\subsection{Version-Bump Retry with Context Reuse}

The workers node implements a three-mode state machine:

\textbf{Mode A (Fresh Start)}: Run all tasks at version 1.

\textbf{Mode B (Retry Only Failed)}: Identify workers with no output file, bump version by 1, run only failed workers, preserve successful outputs.

\textbf{Mode C (Max Retries Exceeded)}: Final attempt, remaining failures get placeholder entries, status set to \texttt{researched\_partial}.

\texttt{CheckpointManager.load\_previous(current\_version)} restores full message history from the previous attempt, resets iteration counter to 0, and injects a context-aware retry message. Constants: \texttt{RESEARCH\_MAX\_VERSIONS=6}, \texttt{RESEARCH\_MAX\_WORKER\_RETRIES=5}, \texttt{WORKER\_TIMEOUT=900}s.

\subsection{Honest Success Reporting}

\texttt{ResearchWorkerRole.parse\_result()} checks \texttt{os.path.isfile()} before reporting success---a binary gate: the worker either has an output file (success) or doesn't (failure). This was the single most impactful reliability fix (v1.4), improving worker success rate from 66.7\% to 100\%.

\subsection{The 5-Dimension Reviewer Scoring System}

The reviewer scores each work on: depth, accuracy, relevance, clarity, originality (each 1--10, max 50 total). A three-tier extraction ensures output: (1) read \texttt{scoring\_report.json}, (2) parse agent messages, (3) programmatic ranking by file existence and summary length with uniform 5/10 scores.

\subsection{LaTeX Report Generation}

The writer uses a sectioned output strategy: each research work is a separate \texttt{.tex} file with \texttt{\textbackslash input\{\}} commands in a main file, avoiding monolithic write truncation. This approach is informed by the agent laboratory paradigm \cite{agentlaboratory} and automated scientific writing pipelines \cite{paperorchestrator}. \texttt{\_latex\_escape()} handles all 10 LaTeX special characters. \texttt{\_fallback\_latex()} generates a complete document from a pre-defined template when LLM generation fails.

\subsection{Always-Forward Routing}

The flow dict encodes a ``progress over perfection'' philosophy: \texttt{researched\_partial $\rightarrow$ reviewer} routes partial worker results to the reviewer rather than aborting. Only three truly unrecoverable states abort the pipeline: no sub-tasks, no reviewable outputs, and no scored works.
